\section{Introduction}
\label{sec:intro}
Reinforcement Learning (RL) has garnered renewed interest in recent
years.  Playing the game of Go \citep{dqn}, robotics control
\citep{robot,ddpg}, and adaptive video streaming \citep{pensieve}
constitute just a few of the vast range of RL applications. Combined
with developments in deep learning, deep reinforcement learning (Deep
RL) has emerged as an accelerator in related fields. From the
well-known success in single-agent deep reinforcement learning, such
as \citet{dqn}, we now witness growing interest in its multi-agent
extension, the multi-agent reinforcement learning (MARL), exemplified
in \citet{gupta2017cooperative,lowe2017multi,coma,omidshafiei2017deep,
  dial, commnet, mordatch2017grounded,havrylov2017,lenient,bicnet,
  lola, tampuu2017multiagent, leibo2017multi,foerster2017stabilising}.
In the MARL problem commonly addressed in these works, multiple agents
interact in a single environment repeatedly and improve their policy
iteratively by learning from observations to achieve a common goal.
Of particular interest is the distinction between two lines of
research: one fostering the direct communication among agents
themselves, as in \citet{dial,commnet} and the other coordinating their
cooperative behavior without direct communication, as in
\citet{foerster2017stabilising,lenient,leibo2017multi}.

In this work, we concern ourselves with the former. We consider MARL
scenarios wherein the task at hand is of a cooperative nature and
agents are situated in a partially observable environment, but each
endowed with different observation power.  We formulate this scenario
into a multi-agent sequential decision-making problem, such that all
agents share the goal of maximizing the same discounted sum of
rewards. For the agents to directly communicate with each other and
behave as a coordinated group rather than merely coexisting
individuals, they must carefully determine the information they
exchange under a practical bandwidth-limited environment and/or in the case
of high-communication cost. To coordinate this exchange of messages,
we adopt the centralized training and distributed execution paradigm
popularized in recent works, {\em e.g.}, \citet{coma, lowe2017multi, vdn,
  qmix, gupta2017cooperative}.

In addition to bandwidth-related constraints, we take the issues of
sharing the communication medium into consideration, especially when
agents communicate over wireless channels. The state-of-the-art standards on wireless communication such as Wi-Fi and LTE specify the way of scheduling users as one of the basic functions. However, as elaborated in
Related work, MARL problems involving scheduling of only a restricted set
of agents have not yet been extensively studied. The key challenges in
this problem are: {\em (i)} that limited bandwidth implies that agents must
exchange succinct information: something concise and yet meaningful
and {\em (ii)} that the shared medium means that potential contenders must
be appropriately arbitrated for proper collision avoidance,
necessitating a certain form of communication scheduling, popularly 
referred to as MAC (Medium Access Control) in the area of wireless
communication.  While stressing the coupled nature of the
encoding/decoding and the scheduling issue, we zero in on the said
communication channel-based concerns and construct our neural network
accordingly.

\myparagraph{Contributions} In this paper, we propose a new deep
multi-agent reinforcement learning architecture, called {\sched}, with
the rationale of centralized training and distributed execution in
order to achieve a common goal better via decentralized cooperation.
During distributed execution, agents are allowed to communicate over
wireless channels where messages are broadcast to all agents in each
agent's communication range. This broadcasting feature of wireless
communication necessitates a Medium Access Control (MAC) protocol to
arbitrate contending communicators in a shared medium. CSMA (Collision
Sense Multiple Access)  in Wi-Fi is one such
MAC protocol. While prior work on MARL to date considers only the
limited bandwidth constraint, we additionally address the shared
medium contention issue in what we believe is the first work of its
kind: which nodes are granted access to the shared medium.
Intuitively, nodes with more important observations should be chosen,
for which we adopt a simple yet powerful mechanism called weight-based
scheduler (WSA), designed to reconcile simplicity in training with
integrity of reflecting real-world MAC protocols in use ({\em e.g.},
802.11 Wi-Fi).  We evaluate {\sched} for two applications: cooperative
communication and navigation and predator/prey and demonstrate that
{\sched} outperforms other baseline mechanisms such as the one without
any communication or with a simple scheduling mechanism such as round
robin.  We comment that {\sched} is not intended for competing with
other algorithms for cooperative multi-agent tasks without considering
scheduling, but a complementary one. We believe that adding our idea
of agent scheduling makes those algorithms much more practical and
valuable.


\myparagraph{Related work} We now discuss the body of relevant
literature. \citet{busoniu2008} and \citet{iql} have studied MARL with
decentralized execution extensively. However, these are based on
tabular methods so that they are restricted to simple environments.
Combined with developments in deep learning, deep MARL algorithms have
emerged \citep{tampuu2017multiagent, coma, lowe2017multi}.
\citet{tampuu2017multiagent} uses a combination of DQN with
independent Q-learning. This independent learning does not perform
well because each agent considers the others as a part of environment
and ignores them. \citet{coma, lowe2017multi, gupta2017cooperative,
  vdn}, and \citet{foerster2017stabilising} adopt the framework of centralized training with
decentralized execution, empowering the agent to learn cooperative
behavior considering other agents' policies without any communication
in distributed execution.

It is widely accepted that communication can further enhance the
collective intelligence of learning agents in their attempt to
complete cooperative tasks. To this end, a number of papers have
previously studied the learning of communication protocols and
languages to use among multiple agents in reinforcement learning. We
explore those bearing the closest resemblance to our research.
\citet{dial,commnet,bicnet, guestrin2002}, and \citet{zhang2013} train multiple
agents to learn a communication protocol, and have shown that
communicating agents achieve better rewards at various tasks.  \citet
{mordatch2017grounded} and \citet{havrylov2017} investigate the possibility of the
artificial emergence of language. Coordinated RL by
\citet{guestrin2002} is an earlier work demonstrating the feasibility
of structured communication and the agents' selection of jointly
optimal action.

Only DIAL \citep{dial} and \citet{zhang2013} explicitly address
bandwidth-related concerns. In DIAL, the communication channel of the
training environment has a limited bandwidth, such that the agents
being trained are urged to establish more resource-efficient
communication protocols. The environment in \citet{zhang2013} also has
a limited-bandwidth channel in effect, due to the large amount of
exchanged information in running a distributed constraint optimization
algorithm. Recently, \citet{attentional_comm} proposes an attentional communication model
that allows some agents who request additional information from others to
gather observation from neighboring agents. However, they
do not explicitly consider the constraints imposed by limited communication bandwidth and/or scheduling due to communication over a shared medium.

To the best of our knowledge, there is no prior work that incorporates an
intelligent scheduling entity in order to facilitate inter-agent
communication in both a limited-bandwidth and shared medium access
scenarios. As outlined in the introduction, intelligent scheduling
among learning agents is pivotal in the orchestration of their
communication to better utilize the limited available
bandwidth as well as in the arbitration of agents contending 
for shared medium access.




%%% Local Variables:
%%% mode: latex
%%% TeX-master: "main"
%%% End:
